\documentclass[full-course-notes.tex]{subfiles}

\begin{document}
\lecturemeta{1}
  {How Programs Run: Compilers, Interpreters, and JIT Compilation (a Java Case Study); the
   Language-Processing System; Why Study Compilers; Kinds of Compilers}
  {CLO 1}
  {\S1.1 (Language Processors), \S1.2 (Language-Processing System -- previewed here, detailed in Lecture 2), \S1.4 (Compiler/Interpreter Distinction), \S1.5 (Applications of Compiler Technology)}
  {[ET] Engineering a Compiler \S1.1; [AP] Modern Compiler Implementation, Ch.\,1}
\lecshort{Compilers, Interpreters \& JIT}
\setcounter{section}{0}
\makelecturetitle

Here is what today actually covers:

\begin{coreidea}[Lecture Roadmap]
\begin{enumerate}
  \item What is a language processor, and what exactly is a \emph{compiler}?
  \item Compilers vs.\ interpreters vs.\ just-in-time (JIT) compilation --- three ways to run a
  program, with a full case study of how \textbf{Java} uses this to become portable.
  \item A \textbf{typical language-processing system}: preprocessor, compiler, assembler,
  linker, and loader -- and exactly what an \textbf{object file} is.
  \item Why study compilers at all, even if you never write one professionally?
  \item Applications of compiler technology far beyond ``turning C into an executable.''
  \item Six different \textbf{kinds of compiler} you'll encounter by name in industry.
\end{enumerate}
This lecture is entirely conceptual --- no algorithms yet. Lecture 2 goes much deeper into the
internal \emph{structure} (phases) of a compiler and the toolchain around it.
\end{coreidea}

\section{What Is a Language Processor?}

\dbtag A \textbf{language processor} is any program that takes another program written in
some source notation and does something useful with it --- runs it, translates it, checks it,
or transforms it. Compilers are the most famous member of this family, but not the only one.

\begin{defbox}{Compiler}
A \textbf{compiler} is a program that reads a program written in one language --- the
\emph{source language} --- and translates it into an equivalent program in another language ---
the \emph{target language}. Along the way, if the source program has errors, the compiler
should report them to the user in a useful way.
\end{defbox}

Three things are packed into that definition and worth pulling apart, because exam questions
love to test exactly this distinction:

\begin{itemize}
  \item \textbf{Translation, not execution.} A compiler's job is to \emph{produce} a target
  program. It does not itself run your code. If you compile \texttt{main.c} with GCC, GCC's job
  ends when it hands you \texttt{a.out}; a completely different program (the operating system's
  loader, then the CPU) is what actually executes it.
  \item \textbf{Equivalence.} The target program must have the same meaning as the source
  program. This is the single hardest engineering promise a compiler makes, and essentially the
  entire field of compiler optimization exists to preserve this promise while still making the
  target program faster.
  \item \textbf{The target language need not be machine code.} Nothing in the definition says
  what the target language \emph{is} -- only that it is ``another language.'' Machine code is
  the single most common and most visible target, which is why people default to thinking
  ``compiler = produces machine code,'' but it is only one case. A compiler whose target is
  \emph{another high-level language} is a source-to-source compiler (\S6); a compiler whose
  target is bytecode for a virtual machine (\texttt{javac} $\to$ JVM bytecode, \S2.3) is still a
  compiler by this same definition, even though no physical CPU can execute its output directly.
\end{itemize}

\begin{center}
\begin{tikzpicture}[
  box/.style={draw, thick, rounded corners, minimum height=1.1cm, minimum width=3.2cm, align=center, fill=nuLight, draw=nuNavy},
  lbl/.style={font=\small\itshape}
]
  \node[box] (src) {Source Program\\\texttt{(e.g.\ C code)}};
  \node[box, fill=dbGreenBg, draw=dbGreen, right=2.6cm of src] (comp) {Compiler};
  \node[box, right=2.6cm of comp] (tgt) {Target Program\\\texttt{(e.g.\ x86 assembly)}};
  \node[below=0.5cm of comp, font=\small, text=examRed] (err) {Error Messages};

  \draw[-{Latex[length=2.5mm]}, thick] (src) -- (comp);
  \draw[-{Latex[length=2.5mm]}, thick] (comp) -- (tgt);
  \draw[-{Latex[length=2.5mm]}, thick, examRed] (comp) -- (err);
\end{tikzpicture}
\end{center}

\begin{examfocus}
Be able to state, in one sentence, that a compiler \emph{translates}; it does not
\emph{execute}. A one-line definition question like ``define a compiler'' is a classic
first-quiz question and students routinely lose marks by describing what an interpreter does
instead -- or by adding ``\ldots into machine code'' to the definition, which is not required
and is contradicted by transpilers and by \texttt{javac} (\S2.3).
\end{examfocus}

\section{Compilers vs.\ Interpreters vs.\ JIT Compilation}

This is the heart of today's lecture (DB \S1.4). There are three broad strategies for getting
from source code to a running program, and almost every real language-processing system you use
day to day is one of these three, or (very commonly) a hybrid.

\subsection{Pure Compilation (Ahead-of-Time)}

\dbtag The compiler translates the entire source program into a target program \emph{before}
execution begins, and the target program can then be run any number of times, independently of
the compiler. In the classic AOT case that gives this strategy its name, the target language is
machine code for the host machine -- this is what C, C++, Rust, and (classic) Fortran compilers
produce -- but the ``ahead-of-time, whole-program-at-once'' strategy itself is independent of
that choice of target (a transpiler, \S6, is also AOT; it just targets another high-level
language instead of machine code).

\begin{center}
\begin{tikzpicture}[
  box/.style={draw, thick, rounded corners, minimum height=1cm, minimum width=2.6cm, align=center},
  phase/.style={box, fill=dbGreenBg, draw=dbGreen},
]
  \node[box, fill=nuLight, draw=nuNavy] (src) {Source Program};
  \node[phase, right=1.8cm of src] (comp) {Compiler};
  \node[box, fill=nuLight, draw=nuNavy, right=1.8cm of comp] (tgt) {Target Program};
  \node[box, fill=nuLight, draw=nuNavy, right=1.8cm of tgt] (input) {Input};
  \node[box, fill=dbGreenBg, draw=dbGreen, below=0.9cm of $(tgt)!0.5!(input)$] (run) {Running the\\Target Program};
  \node[box, fill=nuLight, draw=nuNavy, below=0.9cm of run] (out) {Output};

  \draw[-{Latex[length=2.5mm]}, thick] (src) -- (comp);
  \draw[-{Latex[length=2.5mm]}, thick] (comp) -- (tgt);
  \draw[-{Latex[length=2.5mm]}, thick] (tgt.south) -- ++(0,-0.45) -| (run.north);
  \draw[-{Latex[length=2.5mm]}, thick] (input.south) -- ++(0,-0.45) -| (run.north);
  \draw[-{Latex[length=2.5mm]}, thick] (run) -- (out);
\end{tikzpicture}
\end{center}

\textbf{Advantage -- maximal optimization opportunity.} Because the entire source program is
available before a single instruction runs, an AOT compiler can afford analyses and
transformations that a system under run-time time pressure cannot: multiple, deliberately
ordered optimization passes; \emph{whole-program} (interprocedural) analysis such as global
register allocation or cross-function inlining; and expensive, exhaustive search over code-generation
choices. None of this has to finish before the user sees output, since compilation happens
entirely offline -- only the resulting target program's run time is on the user's clock.

\textbf{Disadvantage -- the output is not portable.} The target program an AOT compiler produces
is machine code for one specific instruction set, calling convention, and (usually) operating
system. A binary built for x86-64 Linux will not run on ARM64 macOS: there is no
platform-independent artifact to ship. Reaching a new target means recompiling from source for
that target (or cross-compiling, \S6) -- unlike an interpreter, which ships the same source to
every platform, or a bytecode/JIT system, which ships the same portable bytecode everywhere.

\textbf{Trade-off, restated:} compilation itself can be slow and happens ``offline,'' but every
run of the target program afterward is fast, because all the translation work -- and all the
optimization it enables -- has already been paid for once, at the cost of tying the result to one
platform.

\subsection{Pure Interpretation}

\dbtag An \textbf{interpreter} does not produce a separate target program at all. Instead, it
directly executes the operations specified in the source program, using the source program (or a
close representation of it) as input on every single run.

\begin{defbox}{Interpreter}
An interpreter is a language processor that reads a source program and, using the input
supplied by the user, \emph{directly} produces the output of that program --- with no separate
translation phase that produces standalone target code the user can run independently.
\end{defbox}

\begin{center}
\begin{tikzpicture}[
  box/.style={draw, thick, rounded corners, minimum height=1cm, minimum width=2.6cm, align=center, fill=nuLight, draw=nuNavy},
  phase/.style={box, fill=dbGreenBg, draw=dbGreen},
]
  \node[box] (src) {Source Program};
  \node[box, below=0.6cm of src] (input) {Input};
  \node[phase, right=2.2cm of $(src)!0.5!(input)$] (interp) {Interpreter};
  \node[box, right=2.2cm of interp] (out) {Output};

  \draw[-{Latex[length=2.5mm]}, thick] (src) -- (interp);
  \draw[-{Latex[length=2.5mm]}, thick] (input) -- (interp);
  \draw[-{Latex[length=2.5mm]}, thick] (interp) -- (out);
\end{tikzpicture}
\end{center}

Classic (older) Python, shell scripts, and simple Basic interpreters are close to this pure
model: the interpreter re-examines and re-dispatches on the source (or its parsed form) every
single time a statement executes, even inside a loop that runs a million times.

\textbf{Advantage -- portability.} The same source program (or bytecode) runs unmodified on any
machine that has a compatible interpreter, with no separate build step per target platform. This
is the mirror image of the AOT compiler's disadvantage above.

\textbf{Disadvantage -- repeated-work overhead.} There is no separate compile step, so you can
start running -- and get feedback on -- code instantly, which is wonderful for interactive
development and REPLs. The cost is that the interpreter re-examines and re-dispatches on the
source (or its parsed form) on \emph{every} execution of a line, with no persistent, whole-program
analysis carried over between runs -- the opposite of an AOT compiler's one-time, multi-pass
optimization. DB notes this repeated-work overhead can make interpretation \textbf{10--100$\times$
slower} than running compiled code.

\begin{examfocus}
Know this speed trade-off in both directions: compilation trades slower ``build time'' for
faster ``run time''; interpretation trades away build time for slower, but more flexible and
immediate, run time. Two further pairings are just as commonly tested: \emph{portability} is a
classic interpreter advantage over AOT compilation (the same source runs unmodified on any
machine that has the interpreter, whereas an AOT binary is tied to one target and needs a
recompile elsewhere), while \emph{optimization opportunity} runs the other way -- an AOT compiler
sees the whole program before run time and can afford multiple, expensive, whole-program passes,
while a pure interpreter re-examines the source on every execution and so cannot afford to.
\end{examfocus}

\subsection{The Hybrid Reality: Bytecode + JIT}

\beyonddb In practice, almost no serious modern language is \emph{purely} compiled or
\emph{purely} interpreted. DB \S1.4 briefly mentions Java's model as a hybrid; it is worth
understanding this hybrid model properly, because it is how the overwhelming majority of code
you run today --- Java, C\#, JavaScript, modern Python (via PyPy), and even Python's own CPython
--- actually executes.

\begin{defbox}{Virtual Machine (VM), vs.\ a Plain Interpreter}
A \textbf{virtual machine}, in the sense used here, is a program that simulates a computer: it
defines its own instruction set (an abstract, software-only ISA) together with the state a real
CPU would have -- a program counter, and either an operand stack (JVM, Wasm) or a set of virtual
registers -- and then runs a fetch-decode-execute loop over that instruction set, exactly the way
real hardware executes machine code, except every step is simulated in software. Crucially, a VM
does \emph{not} interpret your source program directly. There is always a separate front-end
compilation step first (described just below) that translates source into \textbf{bytecode} --
instructions for this abstract machine -- and the VM interprets \emph{that}, not your original
code.

This is different from a ``plain'' (tree-walking) interpreter, of the kind \S2.2 first
introduced: a tree-walking interpreter parses source into an AST and directly, recursively
evaluates that tree (\texttt{evaluate(node)} calling \texttt{evaluate(node.left)}, and so on).
There is no separate instruction set, no fixed opcode encoding, and no simulated CPU state --
the evaluator's own call stack mirrors the shape of the \emph{source program}, not the shape of
a machine. Nothing about it resembles hardware, so it does not earn the name ``virtual machine.''

The distinction, side by side:
\begin{itemize}
  \item \textbf{What is executed.} VM: a fixed, formally-defined bytecode instruction set,
  produced by a separate compile step. Tree-walking interpreter: the source program's own parse
  tree, walked directly, with no separate instruction set at all.
  \item \textbf{Execution state.} VM: an explicit program counter plus an operand stack or
  virtual registers -- state modeled on real hardware. Tree-walker: just the recursive call
  stack of the \texttt{evaluate} function, shaped like the program's grammar, not like a CPU.
  \item \textbf{Is it a genuine compilation target?} VM: yes -- bytecode is often documented and
  shared across multiple front-end languages (JVM bytecode is targeted by Java, Kotlin, Scala,
  and Clojure alike). Tree-walker: no -- the AST is just this one interpreter's private internal
  representation of this one program.
\end{itemize}
Every VM is (or contains) some form of interpreter -- it interprets bytecode instruction by
instruction. But not every interpreter is a VM: a tree-walker interprets source structure directly
and never defines a machine to simulate. (A second, unrelated sense of ``virtual machine'' --
\emph{system} VMs like VMware or QEMU, which virtualize an entire computer including its OS --
is a different concept entirely and not what is meant here.)
\end{defbox}

\begin{enumerate}
  \item \textbf{Source $\to$ bytecode.} A ``front-end'' compiler translates source code into a
  compact, machine-independent intermediate form called \textbf{bytecode} (Java's \texttt{.class}
  files hold JVM bytecode; Python's \texttt{.pyc} files hold Python bytecode).
  \item \textbf{Bytecode $\to$ execution.} A \textbf{virtual machine (VM)} then runs this
  bytecode. It can do this two ways:
  \begin{itemize}
    \item \emph{Interpret} the bytecode directly, instruction by instruction (slow but starts
    instantly) -- this is what most bytecode VMs do at first.
    \item \emph{Just-in-time (JIT) compile} hot pieces of bytecode --- the loops and functions
    that run over and over --- into native machine code \emph{while the program is running}, and
    cache that native code for reuse.
  \end{itemize}
\end{enumerate}

\begin{defbox}{Just-In-Time (JIT) Compilation}
JIT compilation defers translation to native machine code until run time, and does so
selectively: only code that is actually hot (executed frequently) is compiled, while cold code
may remain interpreted. This lets a system get the fast startup of an interpreter and the fast
steady-state speed of a compiler.
\end{defbox}

\begin{coreidea}[Why ``Warm-Up'' Adds a Delay]
The ``time to reach peak speed'' row in \S2.4's table (distinct from startup latency -- see the
exam-focus note there) has a
concrete cause: three costs stack up, in sequence, before any piece of code reaches its fastest
version, and none of them can be skipped.
\begin{enumerate}
  \item \textbf{Early runs pay interpreter-speed cost.} The very first time a method or loop
  executes, no compiled native code for it exists yet -- there is nothing to jump to -- so it
  runs through the interpreter (or a cheap baseline tier, e.g.\ HotSpot's C1), which is
  inherently slower per operation than optimized native code.
  \item \textbf{The system must first observe that code is hot.} A JIT does not compile on the
  first call; it counts. The runtime increments an invocation/back-edge counter on every call or
  loop iteration and only triggers compilation once that counter crosses a threshold. A loop
  that only runs a handful of times may never get compiled at all, because counting to that
  threshold, by construction, requires real executions to have already happened at the slower
  tier.
  \item \textbf{Compilation itself takes real, non-trivial time.} Once flagged hot, the method
  still has to go through the JIT's own pipeline -- build an IR from the bytecode, run
  optimization passes (inlining, register allocation, and more), and emit native code. This
  usually runs on a background thread so it does not block the program outright, but it still
  competes for CPU and memory (see ``Extra runtime overhead'' in \S2.4's table), and the
  \emph{old}, slower version keeps running until the newly compiled version is ready and swapped
  in.
\end{enumerate}
The practical consequence: long-running processes (servers, long batch jobs) amortize this
one-time warm-up cost over a long steady-state run and benefit enormously from JIT compilation,
while short-lived scripts or CLI invocations may finish before ever reaching peak speed --
exactly why AOT-vs.-JIT benchmarks can be misleading if they measure only a program that runs
for a second or two.
\end{coreidea}

\begin{examplebox}{Tiered execution in real systems}
\begin{itemize}
  \item \textbf{JVM (HotSpot):} interprets bytecode first; the C1 (``client'') JIT compiles warm
  methods quickly with light optimization; the C2 (``server'') JIT recompiles the hottest methods
  with heavy optimization once enough profiling data exists.
  \item \textbf{V8 (Chrome/Node.js JavaScript engine):} \texttt{Ignition} interprets bytecode;
  \texttt{Sparkplug} does a fast baseline JIT; \texttt{TurboFan} does an aggressive optimizing
  JIT for hot functions, and can \emph{deoptimize} back to the interpreter if an assumption used
  for optimization (e.g.\ ``this variable is always a number'') turns out to be wrong.
  \item \textbf{CPython (standard Python):} traditionally a pure bytecode interpreter with no
  JIT; \texttt{PyPy} is an alternative Python implementation that adds a JIT and is often
  5--10$\times$ faster on long-running code as a result.
\end{itemize}
\end{examplebox}

\begin{center}
\resizebox{\textwidth}{!}{%
\begin{tikzpicture}[
  box/.style={draw, thick, rounded corners, minimum height=0.95cm, minimum width=2.5cm, align=center, fill=nuLight, draw=nuNavy, font=\small},
  phase/.style={box, fill=dbGreenBg, draw=dbGreen},
  hot/.style={box, fill=examRedBg, draw=examRed},
]
  \node[box] (src) {Source\\Program};
  \node[phase, right=1.6cm of src] (fe) {Front-end\\Compiler};
  \node[box, right=1.6cm of fe] (bc) {Bytecode};
  \node[phase, below right=1.1cm and 0.0cm of bc] (vm) {Bytecode\\Interpreter};
  \node[hot, right=1.6cm of bc] (jit) {JIT\\Compiler};
  \node[box, right=1.6cm of jit] (native) {Native\\Machine Code};

  \draw[-{Latex[length=2.5mm]}, thick] (src) -- (fe);
  \draw[-{Latex[length=2.5mm]}, thick] (fe) -- (bc);
  \draw[-{Latex[length=2.5mm]}, thick] (bc) -- (jit);
  \draw[-{Latex[length=2.5mm]}, thick] (jit) -- (native);
  \draw[-{Latex[length=2.5mm]}, thick] (bc.south) -- (vm.north);
  \node[below=0.15cm of vm, font=\scriptsize\itshape, text width=2.6cm, align=center] {runs cold /\\one-shot code};
  \node[below=0.15cm of native, font=\scriptsize\itshape, text width=2.6cm, align=center] {runs hot,\\repeated code};
\end{tikzpicture}}
\end{center}

\begin{enrichment}[Extra Depth: AOT with a JIT Fallback]
DB draws the compiler/interpreter line fairly simply. In practice, the line is blurrier still:
some modern systems (GraalVM Native Image, .NET's ReadyToRun) precompile most code AOT for fast
startup but still keep a JIT available for code discovered or specialized at run time. So
``compiled'' and ``JIT-compiled'' are not always mutually exclusive labels for one system.
\emph{(Transpilers, cross-compilers, bootstrapping compilers, and decompilers are their own
category of ``kind of compiler'' rather than a compiler/interpreter distinction -- we cover all
of these properly later in this lecture.)}
\end{enrichment}

\subsection{Comparing the Three Models}

\small
\renewcommand{\arraystretch}{1.3}
\begin{longtable}{@{} p{0.16\textwidth} p{0.27\textwidth} p{0.27\textwidth} p{0.27\textwidth} @{}}
\caption{Comparing the three execution models} \label{tab:compare-three-models} \\
\toprule
 & \textbf{Pure Compiler} & \textbf{Pure Interpreter} & \textbf{JIT / Hybrid} \\
\midrule
\endfirsthead

\multicolumn{4}{@{}l}{\itshape Table \thetable{} -- comparing the three execution models (continued)} \\
\toprule
 & \textbf{Pure Compiler} & \textbf{Pure Interpreter} & \textbf{JIT / Hybrid} \\
\midrule
\endhead

\midrule
\multicolumn{4}{@{}r}{\itshape continued on next page} \\
\endfoot

\bottomrule
\endlastfoot

\textbf{When translation happens} & Entirely before run time & Continuously, during run time & Source $\to$ bytecode happens before run time; bytecode $\to$ native code (the JIT step itself) is \textbf{always fully during} run time -- see note below \\
\textbf{Time to reach peak speed} & N/A -- already at peak from instruction one (see above) & N/A -- there is no ``peak''; every instruction runs at the same interpreted speed forever & \textbf{Nonzero (``warm-up'')} -- profiling must first identify hot code, then compilation must finish, before that code is upgraded to native speed; the program is already running the whole time this happens \\
\textbf{Steady-state speed} & Fastest & Slowest (re-dispatch overhead) & Approaches compiled speed for hot code \\
\textbf{Portability} & Target is machine-specific -- a new platform needs a recompile (or cross-compile) & Source runs anywhere the interpreter exists & Bytecode is portable; the JIT-generated native code is not, but it never has to be shipped \\
\textbf{Optimization opportunity} & Largest -- whole program visible before run time, so multiple, expensive, whole-program passes are affordable & Smallest -- source (or bytecode) is re-examined on every execution, so heavyweight analysis is normally too costly to repeat & Bounded by run-time budget -- fewer/cheaper passes than AOT and limited mostly to hot methods, but compensated by real profiling data (next row) unavailable to AOT \\
\textbf{Can use runtime profiling info?} & No -- only sees the source, never the running program & N/A -- re-examines source every time, no persistent profile & \textbf{Yes} -- observes actual branch frequencies and value types, enabling speculative optimizations no static compiler can safely make \\
\textbf{Extra runtime overhead} & None -- compiler is long gone by run time & Dispatch overhead only & Profiling counters, background compiler threads, and a generated-code cache all compete with the program for CPU/memory \\
\textbf{Debuggability / interactivity} & Weakest (must recompile to test a change) & Best (REPL, immediate feedback) & Good (often still supports a REPL over bytecode) \\
\textbf{Typical users} & C, C++, Rust, Fortran & Shell scripts, classic BASIC & Java, C\#, JavaScript, PyPy \\
\end{longtable}
\normalsize

\begin{examfocus}
The ``When translation happens'' row is the row most often misread. The hybrid pipeline has
\emph{two} separate translation stages, and only one of them is ``just-in-time'':
\begin{itemize}
  \item \textbf{Source $\to$ bytecode} (e.g.\ \texttt{javac} producing \texttt{.class} files) is
  ordinary ahead-of-time compilation. It happens once, before the program is ever run, exactly
  like \S2.1's pure-compiler case -- it is just that its \emph{target} language is bytecode
  instead of machine code.
  \item \textbf{Bytecode $\to$ native code}, the actual JIT step, is \emph{never} done ahead of
  time, not even partially. The decision of \emph{which} loop or method is hot enough to be
  worth compiling is made by watching the program actually run (via profiling counters), and the
  native code generated is specific to the behavior observed \emph{during that run} -- this is
  exactly what ``just-in-time'' means: translation deferred until the last possible moment,
  triggered by run-time evidence, not decided in advance.
\end{itemize}
So the ``partly before, partly during'' language in the table describes the \emph{full hybrid
pipeline} (front-end compile, then JIT), not the JIT step in isolation. Taken alone, JIT
compilation is entirely a run-time activity, no different in \emph{when} it happens from pure
interpretation -- what makes it different is \emph{what} it produces (a reusable native-code
artifact) rather than \emph{when} it produces it.
\end{examfocus}

\begin{examfocus}
A related trap, worth knowing even though the table above does not track it separately:
``startup latency'' and ``warm-up delay'' sound like the same thing but are not. \textbf{Startup
latency} asks only: how long before the program starts executing at all? By that measure, JIT
and pure interpretation are identical -- a JIT-based system \emph{begins} by interpreting (or
running a cheap baseline tier) exactly like a pure interpreter does, so there is no extra pause
before the first instruction runs. ``Warm-up,'' by contrast, is not a delay \emph{before}
starting; it is the time \emph{after} the program has already started before some of its code
gets promoted to native speed -- exactly what the ``Time to reach peak speed'' row above
captures. The program is running the entire time warm-up is happening -- it is simply running at
interpreted (not yet peak) speed for that stretch. So it is correct to say a JIT system has low
startup latency \emph{and} a nonzero warm-up period -- these are not in tension, because they
answer two different questions.
\end{examfocus}

\begin{historynote}
Interpretation is not a modern shortcut invented for scripting languages --- it is \emph{older}
than compilation. Early LISP (1958) was interpreted because writing a LISP compiler was itself a
research problem; interpreters let you experiment with a language before anyone knew how to
compile it well. JIT compilation is also older than most people assume: the term traces to
Smalltalk and self-modifying-code systems of the 1980s, and was popularized for the mainstream by
Sun's HotSpot JVM in 1999 -- it was Java's answer to ``we want C-like speed but Smalltalk-like
portability and safety.''
\end{historynote}

\subsection{Case Study: How Java Achieves Portability}

\beyonddb Java is the single clearest real-world illustration of everything in \S2, so it is
worth walking through as a complete case study --- especially since Java's own design goals were
\emph{explicitly} about the compiler/interpreter trade-off from the start.

\subsubsection{From Pure Interpreter to JIT: A Short History}

\begin{center}
\renewcommand{\arraystretch}{1.35}
\begin{tabularx}{\textwidth}{@{}l X@{}}
\toprule
\textbf{Era} & \textbf{What Happened} \\
\midrule
Java 1.0 (1996) & The original JVM was a \textbf{pure bytecode interpreter}. Every bytecode
instruction was fetched, decoded, and dispatched on every single execution -- simple and
portable, but roughly 10--20$\times$ slower than equivalent compiled C, which badly hurt Java's
early reputation for performance. \\
Java 1.2 -- ``HotSpot'' research & Sun's research team built an adaptive compiler codenamed
\textbf{HotSpot}: instead of compiling everything, or nothing, it would \emph{watch the program
run} and compile only the small fraction of methods that actually matter for performance. \\
Java 1.3 (2000) onward & HotSpot became the default production JVM. This is the point Java
became a genuine \textbf{hybrid}: start by interpreting (fast startup), then JIT-compile
whichever methods turn out to be ``hot'' (fast steady state). Every mainstream JVM since has
kept this same architecture. \\
\bottomrule
\end{tabularx}
\end{center}

\subsubsection{Compilation to Bytecode, Then Portability via the JVM}

\dbtag \texttt{javac}, the Java \emph{compiler}, does not produce machine code for any real CPU
at all. It produces \textbf{Java bytecode} -- instructions for an abstract, idealized machine
called the \textbf{Java Virtual Machine (JVM)} that does not physically exist as hardware.

\begin{center}
\resizebox{\textwidth}{!}{%
\begin{tikzpicture}[
  box/.style={draw, thick, rounded corners, minimum height=1cm, minimum width=3cm, align=center, font=\small},
  tool/.style={box, fill=dbGreenBg, draw=dbGreen},
  data/.style={box, fill=nuLight, draw=nuNavy},
  hot/.style={box, fill=examRedBg, draw=examRed},
  >=Latex
]
  \node[data] (src) at (0,0) {\texttt{Hello.java}\\(source)};
  \node[tool] (javac) at (4,0) {\texttt{javac}\\(compiler)};
  \node[data] (class) at (8,0) {\texttt{Hello.class}\\(JVM bytecode)};
  \node[tool] (interp) at (12.3,1.1) {JVM:\\Interpreter};
  \node[hot]  (jit)    at (12.3,-1.1) {JVM:\\JIT Compiler};
  \node[data] (out) at (16.5,0) {Running on\\\emph{this} CPU};

  \draw[->, thick] (src) -- (javac);
  \draw[->, thick] (javac) -- (class);
  \draw[->, thick] (class.east) -- ++(0.7,0) |- (interp.west);
  \draw[->, thick] (class.east) -- ++(0.7,0) |- (jit.west);
  \draw[->, thick] (interp.east) -| (out.north);
  \draw[->, thick] (jit.east) -| (out.south);
\end{tikzpicture}}
\end{center}

\begin{coreidea}[``Write Once, Run Anywhere'']
Because \texttt{javac} only ever has to target \emph{one} abstract machine (the JVM
specification) instead of every physical CPU family, the exact same \texttt{Hello.class} file
runs unmodified on Windows/x86, Linux/ARM, or a JVM on any other platform --- as long as
\emph{that} platform has a compliant JVM installed. All the platform-specific work is pushed
into the JVM implementation itself (one per platform, built once by the JVM vendor), instead of
into every single Java program ever written. This is Sun Microsystems' famous ``Write Once, Run
Anywhere'' pitch for Java, and it is the exact same $M + N$ argument (rather than $M \times N$)
you will meet again with LLVM IR in Lecture 2 --- just applied to distribution instead of to
compiler internals.
\end{coreidea}

\subsubsection{HotSpot's Tiered JIT in More Depth}

\dbtag Modern HotSpot uses \textbf{tiered compilation}, layering the interpreter and two JIT
compilers so the JVM is never stuck doing more work than the code's importance justifies:

\begin{enumerate}
  \item Every method starts out \textbf{interpreted}, while the JVM silently maintains an
  \textbf{invocation counter} (calls) and a \textbf{back-edge counter} (loop iterations) for it.
  \item Once a counter crosses a threshold, the \textbf{C1 (client) compiler} quickly produces a
  lightly-optimized native version, and also inserts \textbf{profiling instrumentation} into it
  (recording which branches are taken, which concrete types show up at a call site, \dots).
  \item If the method keeps getting hotter, the \textbf{C2 (server) compiler} recompiles it a
  second time, now using the profiling data gathered by the C1 version to make aggressive,
  \emph{speculative} optimizations: inlining a virtual method call that has -- so far -- always
  resolved to the same concrete type; eliminating a heap allocation entirely if it can prove the
  object never escapes the method (\textbf{escape analysis}); unrolling a hot loop.
  \item If a speculative assumption is later violated (a second, different type shows up at that
  call site), HotSpot \textbf{deoptimizes}: it discards the C2 code for that method and falls
  back to the interpreter (or C1) until it is confident enough to try recompiling again.
\end{enumerate}

\begin{examfocus}
Be able to explain, in your own words, why HotSpot's C2 compiler can safely make optimizations
(like assuming a call site's type never changes) that a purely static AOT compiler like GCC
never could -- and what mechanism (deoptimization) lets it recover if that assumption turns out
to be wrong.
\end{examfocus}

\subsubsection{JIT vs.\ AOT vs.\ Interpreter: Advantages and Disadvantages}

Pulling together everything from the enhanced table in \S2.4 and the HotSpot walkthrough above,
here is the direct trade-off summary for a JIT compiler specifically:

\begin{center}
\renewcommand{\arraystretch}{1.3}
\begin{tabularx}{\textwidth}{@{}X X@{}}
\toprule
\textbf{JIT Advantages} & \textbf{JIT Disadvantages} \\
\midrule
Can use \textbf{real runtime profile data} (actual types, actual branch frequencies) to make
optimizations no static compiler can safely attempt -- generally impossible for an AOT compiler,
which only ever sees the source. &
\textbf{Warm-up latency}: peak performance is only reached after enough executions to get hot
code recompiled -- a real problem for short-lived programs (e.g.\ CLI tools, serverless
functions). \\
Ships one \textbf{portable} bytecode artifact instead of a separate binary per target CPU --
matching an interpreter's portability, unlike AOT. &
Consumes \textbf{extra CPU and memory at run time} for profiling counters, the JIT compiler
threads themselves, and a generated-code cache -- resources an AOT binary or a plain interpreter
never spend. \\
Reaches \textbf{steady-state speed close to AOT-compiled native code} for long-running,
hot-path-dominated programs -- far faster than a pure interpreter ever gets. &
\textbf{Less predictable} performance moment-to-moment (a method might pause execution while it
gets recompiled), which matters for hard real-time systems that AOT compilation suits better. \\
Can \textbf{adapt to the actual host CPU} at run time (e.g.\ use available vector instructions)
without shipping separate binaries per microarchitecture, unlike a single AOT build. &
\textbf{Significantly more implementation complexity}: an interpreter, multiple JIT tiers,
profiling instrumentation, and deoptimization machinery, versus AOT's comparatively simple
``compile once, done.'' \\
\bottomrule
\end{tabularx}
\end{center}

\begin{enrichment}[Extra Depth: The Java Class File, Briefly]
A \texttt{.class} file is itself a small, well-defined binary format worth knowing at a glance:
it opens with a \textbf{magic number} (\texttt{0xCAFEBABE}, a Sun Microsystems in-joke), a
class-file version number, a \textbf{constant pool} (a table holding every literal, class name,
and method/field reference the class uses -- similar in spirit to the symbol table we will meet
in Lecture 2), and then the actual bytecode for each method. Tools like \texttt{javap -c
Hello.class} let you disassemble a compiled class and read its raw JVM bytecode instructions
directly -- a nice hands-on way to see \S2's ``token stream $\to$ bytecode'' pipeline made
completely concrete.
\end{enrichment}

\section{Why Study Compilers? (DB \S1.1)}

\dbtag You may never write a production C compiler. So why is this course core, and why does
almost every strong CS program require it? Three overlapping reasons:

\begin{enumerate}
  \item \textbf{Compilers are, in DB's phrase, ``one of the largest and most complex software
  systems'' commonly built} -- yet a working one is buildable in a single semester project.
  Nowhere else in the curriculum will you build something this large, this precisely specified,
  end to end, from a formal problem statement to a working artifact.
  \item \textbf{The techniques generalize far beyond compilers.} Finite automata power text
  editors, network protocol parsers, and input validation. Context-free grammars appear in JSON
  parsers, config-file readers, and query languages. Data-flow analysis underlies static
  analyzers, IDE ``find usages'' features, and security vulnerability scanners. You are not just
  learning to build a compiler; you are learning the standard toolkit for building \emph{any}
  system that must process structured, language-like input reliably.
  \item \textbf{Compilers sit at the hardware/software boundary.} Understanding how a
  \texttt{for} loop becomes register moves and branch instructions demystifies performance --
  why one piece of ``equivalent'' code runs 10$\times$ faster than another, why cache-friendly
  code matters, and why some languages can be faster than others for reasons that have nothing to
  do with the programmer's skill.
\end{enumerate}

\begin{enrichment}[Extra Depth: A ``Meta'' Argument]
There's a fourth reason DB does not spell out explicitly but that experienced engineers often
cite: compiler construction is one of the few courses where you are forced to build a
\emph{multi-stage pipeline} with cleanly separated concerns (lexing, parsing, semantic
analysis, code generation) and \emph{formally specified} correctness at each stage. This is
directly transferable to designing any large pipeline system -- data pipelines, network stacks,
even ML training pipelines -- where the same discipline of ``define an intermediate
representation, verify it, transform it, repeat'' shows up again and again.
\end{enrichment}

\section{Applications of Compiler Technology (DB \S1.5)}

\dbtag DB is explicit that compiler techniques reach ``far beyond'' translating a
general-purpose programming language into machine code. Five application areas DB highlights:

\begin{enumerate}
  \item \textbf{Implementation of high-level programming languages.} The obvious case, but note
  the historical point DB makes: higher-level languages raised the level of abstraction
  precisely \emph{because} compiler technology matured enough to make that abstraction
  affordable in performance terms.
  \item \textbf{Optimizations for computer architectures.} As hardware has grown more parallel
  (multi-core, SIMD/vector units, deep pipelines, GPUs), the compiler has taken on more
  responsibility for finding parallelism and hiding memory latency that programmers cannot
  realistically manage by hand.
  \item \textbf{Design of new computer architectures.} Hardware designers now co-design
  instruction sets \emph{with} the compiler in mind (RISC philosophy: keep hardware simple, let
  the compiler do the scheduling), since a feature only helps if a compiler can reliably exploit
  it.
  \item \textbf{Program translations.} Beyond source-to-machine-code, DB lists: binary
  translation (retargeting compiled binaries to a new instruction set, e.g.\ Apple's Rosetta 2
  translating x86 binaries to run on Apple Silicon), hardware synthesis (compiling a hardware
  description language like Verilog into actual circuit layouts), database query interpreters
  (translating SQL into an execution plan), and compiled simulation (compiling a simulation
  model, e.g.\ for chip verification, into fast executable code rather than interpreting it).
  \item \textbf{Software productivity tools.} Type checking, bounds checking, and static
  program-analysis tools that catch bugs before run time all reuse the same front-end machinery
  (lexers, parsers, semantic analyzers) that compilers use.
\end{enumerate}

\begin{enrichment}[Extra Depth: Applications DB Does Not Mention]
The five areas above were written with early-2000s computing in mind. Several major
compiler-technology application areas have grown enormously since and deserve a mention:
\begin{itemize}
  \item \textbf{Machine-learning compilers.} Frameworks like Google's XLA, Apache TVM, and
  MLIR (Multi-Level Intermediate Representation, itself built using classic compiler theory)
  compile neural-network computation graphs down to highly optimized GPU/TPU kernels -- this is
  now one of the most active areas of compiler research.
  \item \textbf{WebAssembly (WASM).} A portable, sandboxed bytecode target designed so
  \emph{any} source language (C, Rust, Go, \dots) can compile to it and run safely inside a
  browser or standalone runtime at near-native speed.
  \item \textbf{Superoptimizers.} Tools (e.g.\ Souper, STOKE) that search -- sometimes with
  SMT solvers, sometimes with reinforcement learning -- for provably optimal short instruction
  sequences, going beyond what traditional rule-based optimization passes can find.
  \item \textbf{Domain-specific languages (DSLs) and their compilers.} Halide (image-processing
  pipelines), SQL query planners, regular-expression engines, and shader languages (GLSL/HLSL for
  GPUs) are all compilers for a language deliberately kept small and specialized so it can be
  optimized far more aggressively than a general-purpose language could be.
  \item \textbf{Security applications.} Control-flow integrity instrumentation, automatic
  bounds-check insertion (AddressSanitizer), and fuzzing harness generation are compiler-pass
  techniques applied specifically to hardening software against attack.
  \item \textbf{LLVM as shared infrastructure.} Perhaps the biggest structural change since DB
  2nd edition: LLVM lets dozens of front-ends (Clang/C/C++, Rust, Swift, Julia, and more) share
  one very well-engineered optimizer and set of back-ends, instead of every language reinventing
  code generation. We will return to LLVM in Lecture 25.
\end{itemize}
\end{enrichment}

\section{A Typical Language-Processing System}

\dbtag The ``compiler'' most people invoke by typing \texttt{gcc file.c} is not, strictly, one
single program at all -- it is a \textbf{driver} that coordinates several distinct tools in
sequence. DB places the compiler proper inside this larger \textbf{language-processing system}.

\begin{center}
\resizebox{\textwidth}{!}{%
\begin{tikzpicture}[
  box/.style={draw, thick, rounded corners, minimum height=1.1cm, minimum width=2.7cm, align=center, font=\small},
  tool/.style={box, fill=dbGreenBg, draw=dbGreen},
  data/.style={box, fill=nuLight, draw=nuNavy},
  >=Latex
]
  \node[data] (src)  at (0,0)  {Source Program\\(with macros,\\\texttt{\#include}s)};
  \node[tool] (pre)  at (4,0)  {Preprocessor};
  \node[data] (src2) at (8,0)  {Pure Source\\Program};
  \node[tool] (comp) at (12,0) {Compiler};
  \node[data] (asm)  at (16,0) {Target\\Assembly Program};

  \node[tool] (assm) at (12,-2.4) {Assembler};
  \node[data] (rel)  at (16,-2.4) {Relocatable\\Object Code};
  \node[data] (libs) at (8,-4.6)  {Library Files \&\\Relocatable Object Files};
  \node[tool] (ll)   at (12,-4.6) {Linker / Loader};
  \node[data] (tgt)  at (16,-4.6) {Target\\Machine Code};

  \draw[->, thick] (src) -- (pre);
  \draw[->, thick] (pre) -- (src2);
  \draw[->, thick] (src2) -- (comp);
  \draw[->, thick] (comp) -- (asm);
  \draw[->, thick] (asm.south) -- ++(0,-0.5) -| (assm.north);
  \draw[->, thick] (assm) -- (rel);
  \draw[->, thick] (rel.south) -- ++(0,-0.5) -| (ll.north);
  \draw[->, thick] (libs) -- (ll);
  \draw[->, thick] (ll) -- (tgt);
\end{tikzpicture}}
\end{center}
\begin{center}
\small\textit{A typical language-processing system (compare DB Fig.\ 1.5): the ``compiler'' you
invoke on the command line is really a pipeline of five distinct programs.}
\end{center}

\begin{itemize}
  \item \textbf{Preprocessor.} Expands macros (\texttt{\#define}), includes files
  (\texttt{\#include}), and resolves conditional-compilation directives (\texttt{\#ifdef}) ---
  all as pure text substitution, before the compiler's grammar-aware phases ever run.
  \item \textbf{Compiler (proper).} Everything from \S1: source $\to$ target assembly (today's
  ``black box''; Lecture 2 opens it up phase by phase).
  \item \textbf{Assembler.} Translates human-readable assembly instructions into
  \textbf{relocatable machine code} and packages it as an \textbf{object file}.
  \item \textbf{Linker.} Combines multiple object files and library files into one program,
  resolving references between them.
  \item \textbf{Loader.} Places the final program into memory and starts it running.
\end{itemize}

\begin{enrichment}[Is an Assembler a Compiler?]
By the definition in \S1, yes: an assembler reads a program in one language (assembly) and
translates it into an equivalent program in another language (relocatable machine code), with no
execution of its own along the way -- that is exactly the compiler definition. It is given its
own name and its own box in the pipeline diagram above for practical, not definitional, reasons:
assembly is designed to be an almost line-for-line, one-to-one notation for machine instructions,
so an assembler's job skips most of what makes a ``full'' compiler hard (there is essentially no
analysis phase, no optimization search, and little structural gap between source and target to
bridge) and reduces mostly to table lookup and address bookkeeping. The lesson generalizes:
\emph{compiler} names a \emph{role} -- translate one language into another, without executing
either -- not a claim about how sophisticated the translation has to be. An assembler, a
JVM-bytecode compiler like \texttt{javac}, and a heavily-optimizing tool like GCC are all
compilers in exactly the same technical sense; they differ only in how much analysis and
transformation the translation requires.
\end{enrichment}

\begin{defbox}{Object File}
An \textbf{object file} is the output of the assembler: a binary file containing machine
instructions and data \emph{that is not yet a runnable program}. Critically, it contains:
\begin{itemize}
  \item \textbf{Relocatable code/data}, with memory addresses left as placeholders, since the
  assembler does not know where in memory this code will ultimately live;
  \item A table of \textbf{symbols this file defines} (e.g.\ a function it implements, available
  for other files to call); and
  \item A table of \textbf{symbols this file references but does not define} (e.g.\ a call to
  \texttt{printf}, defined elsewhere in the C standard library) -- these are exactly what the
  linker must resolve next.
\end{itemize}
On Linux this is the \texttt{.o} file produced by \texttt{gcc -c}; the same idea is called an
\texttt{.obj} file on Windows.
\end{defbox}

\begin{coreidea}[What the Linker Actually Does]
Programs are almost never compiled as a single file. The \textbf{linker}'s job is to take
\emph{several} object files (and pre-built library files) and merge them into one program, by:
\begin{enumerate}
  \item \textbf{Symbol resolution:} matching every symbol a file \emph{references} (like a call
  to a function defined in another file, or in a library) to the file that actually
  \emph{defines} it. An unresolved symbol here is exactly the ``undefined reference to
  \texttt{foo}'' error you see at link time, as distinct from a compiler error.
  \item \textbf{Relocation:} now that every piece of code has been assigned a real, final
  address in the combined program, the linker patches every placeholder address left by the
  assembler with the correct final value.
\end{enumerate}
The output is one combined file --- an \textbf{executable} --- with (in the simplest case) no
unresolved symbols and no remaining placeholder addresses.
\end{coreidea}

\begin{coreidea}[What the Loader Actually Does]
The \textbf{loader} is typically part of the operating system, not the compiler toolchain. When
you run a program, the loader:
\begin{enumerate}
  \item Allocates memory for the program (code, data, stack, heap regions);
  \item Copies (or memory-maps) the executable's instructions and initial data into that memory;
  \item Performs any \emph{final} address fix-ups that were deferred until this exact moment
  (relevant for dynamically linked libraries, previewed in Lecture 2); and
  \item Transfers control to the program's entry point -- the actual first instruction that
  runs, which is \emph{not} always your \texttt{main} function directly (a small
  language-runtime startup routine usually runs first).
\end{enumerate}
\end{coreidea}

\begin{examfocus}
A very common exam trap: ``which phase reports an `undefined reference' / `unresolved external
symbol' error for a missing function?'' The answer is the \textbf{linker}, not the compiler ---
the compiler is perfectly happy as long as a function is \emph{declared} (e.g.\ via a header
file); only the linker checks that a matching \emph{definition} actually exists somewhere among
all the object files and libraries it was given.
\end{examfocus}

\begin{enrichment}[A Preview of What's Coming in Lecture 2]
This is deliberately the \emph{introductory} pass over the toolchain. Lecture 2 goes much
further: static vs.\ dynamic linking, position-independent code and ASLR, real object-file
formats (ELF, Mach-O, PE), and how build systems like Make, Ninja, and Bazel orchestrate this
whole pipeline across thousands of files.
\end{enrichment}

\section{Different Kinds of Compilers}

\beyonddb ``Compiler'' is used as an umbrella term across the industry for several genuinely
different tools, distinguished by \emph{what kind of language} sits on each side of the
translation (high-level vs.\ low-level) and \emph{which machine} runs the result. Six worth
knowing by name:

\begin{center}
\renewcommand{\arraystretch}{1.45}
\begin{tabularx}{\textwidth}{@{}p{2.7cm} p{3.6cm} X@{}}
\toprule
\textbf{Kind} & \textbf{Source $\to$ Target} & \textbf{Example(s)} \\
\midrule
\textbf{Native (AOT) compiler} & High-level $\to$ machine code, same target as the host machine
& GCC, Clang, \texttt{rustc} compiling for the machine they run on \\
\textbf{Source-to-source compiler} (\textbf{transpiler}) & High-level $\to$ a \emph{different}
high-level language & TypeScript $\to$ JavaScript; Babel (modern JS $\to$ older JS); CoffeeScript
$\to$ JavaScript \\
\textbf{Cross-compiler} & High-level $\to$ machine code for a target \emph{different} from the
host running the compiler & Compiling ARM binaries for a phone on an x86 laptop; embedded
firmware toolchains \\
\textbf{Binary translator} & Low-level (one instruction set) $\to$ low-level (a \emph{different}
instruction set) & Apple's Rosetta 2 (x86-64 $\to$ ARM64); QEMU's user-mode emulation \\
\textbf{Decompiler} & Low-level (machine code or bytecode) $\to$ high-level-\emph{like} source &
Ghidra, IDA Pro (machine code); JD-GUI, CFR (Java bytecode) \\
\textbf{Bootstrapping (self-hosting) compiler} & A compiler for language X, \emph{written in} X
itself & GCC (mostly C); \texttt{rustc} (written in Rust) \\
\bottomrule
\end{tabularx}
\end{center}

\begin{examplebox}{Two of These, Worked Through}
\textbf{Binary translation (Rosetta 2):} when Apple moved Macs from Intel (x86-64) to Apple
Silicon (ARM64) in 2020, millions of existing apps were compiled x86-64 binaries with no source
available to the end user. Rosetta 2 translates those already-compiled x86-64 binaries into
ARM64 machine code -- notice both sides of this translation are \emph{low-level}, which is what
makes it a binary translator rather than an ordinary compiler.

\textbf{Bootstrapping (GCC):} you cannot compile the very first C compiler \emph{using} a C
compiler, since none exists yet. Historically, an initial, minimal C compiler was written in
assembly or an even earlier language; once that minimal compiler could compile \emph{a subset}
of C, a fuller C compiler was written in C and compiled by the minimal one; that fuller compiler
then compiles \emph{itself} (and every subsequent version) from then on. This
chicken-and-egg-breaking process is exactly what ``self-hosting'' means.
\end{examplebox}

\begin{examfocus}
Given a described tool (``translates Kotlin to Java source,'' ``translates ARM machine code
back into readable pseudo-C,'' ``compiles Swift for an iPhone from a Mac''), be able to name
which of these six kinds it is and justify your answer using the \emph{source $\to$ target}
column above -- this classification-by-example is a very common short-answer format.
\end{examfocus}

\begin{enrichment}[Where the Compiler-vs-Interpreter and Kind-of-Compiler Axes Meet]
These two classification schemes are independent and can combine freely. A JIT compiler
\emph{is} a native compiler by this section's scheme (high-level/bytecode $\to$ machine code),
just one whose \emph{timing} (during execution) differs from an AOT native compiler's. A
transpiler like Babel is almost always itself an AOT tool -- it finishes its work entirely
before the JavaScript it produces ever runs. Don't confuse ``what kind of translation is this''
(this section) with ``when does the translation happen'' (\S2) -- exam questions sometimes test
both axes on the same tool deliberately.
\end{enrichment}

\section{A Preview: What's Next}

Lecture 2 opens up the compiler ``black box'' from \S1 of this lecture and walks through its
internal phases -- lexical analysis, parsing, semantic analysis, intermediate code generation,
optimization, and final code generation -- in real depth. It also returns to the
language-processing toolchain from \S6, going far beyond today's introduction: static vs.\
dynamic linking, position-independent code, real object-file formats, and how build systems
orchestrate compilation across large projects. Everything in today's lecture is about
\emph{what} these systems do from the outside; Lecture 2 begins answering \emph{how}, in detail.

\section{Self-Check Questions}

\begin{enrichment}[Practice -- Not Graded]
\begin{enumerate}
  \item In one sentence, explain why ``a compiler executes your program'' is a factually
  incorrect statement, and what does execute it instead.
  \item Classic CPython has no JIT; PyPy does. Predict, with a reason, which one would be faster
  for a program consisting of a single loop that runs 50 million times, versus a script that
  runs once and exits immediately.
  \item Explain, using the $M+N$ vs.\ $M \times N$ argument, why compiling to JVM bytecode makes
  Java portable in a way that compiling directly to x86 machine code would not.
  \item A program compiles successfully with no errors, but fails at link time with ``undefined
  reference to \texttt{compute\_total}.'' Which tool caught this, and why couldn't the compiler
  have caught it earlier?
  \item Is a transpiler (e.g.\ TypeScript $\to$ JavaScript) a compiler under DB's definition?
  Justify your answer using the definition given in \S1, and separately, name which of the six
  ``kinds of compiler'' from \S6 it belongs to.
  \item Give one example, not listed in this lecture, of a real tool on your own laptop that is
  best described as an ``interpreter'' rather than a compiler.
  \item HotSpot's C2 JIT only compiles a method after it has been \emph{interpreted} many
  times. Why not just JIT-compile every method the very first time it's called?
\end{enumerate}
\end{enrichment}

\begin{takeaways}
\begin{itemize}
  \item A \textbf{compiler translates}; it does not itself execute the source program. An
  \textbf{interpreter} executes the source directly, with no independent target program left
  behind.
  \item Compilers trade slower build time for faster, repeatable run time, at the cost of a
  target binary that is \textbf{not portable} (tied to one ISA/OS) and paid for by the whole-program
  visibility that enables aggressive, multi-pass optimization. Interpreters trade that away for
  instant startup, easy interactivity, and \textbf{source-level portability}, at the cost of
  repeated re-interpretation overhead on every run.
  \item Most real systems today (JVM, V8, .NET, PyPy) are \textbf{hybrids}: compile to portable
  bytecode, interpret it at first, then \textbf{JIT-compile} the hot paths to native code at run
  time -- Java's own history (pure interpreter in 1.0, HotSpot's tiered JIT from 1.3 onward) is
  the clearest real example, and it's what makes ``write once, run anywhere'' work.
  \item A JIT's big edge over AOT is \textbf{using real runtime profile data} for speculative
  optimization; its costs are \textbf{warm-up latency} and \textbf{runtime profiling/compilation
  overhead} -- know this trade-off table cold.
  \item The compiler proper sits inside a larger \textbf{language-processing system}:
  preprocessor $\to$ compiler $\to$ assembler $\to$ linker $\to$ loader. The assembler produces
  an \textbf{object file} (relocatable code plus defined/referenced symbol tables); the
  \textbf{linker} resolves symbols and relocates addresses; the \textbf{loader} places the
  program in memory and starts it.
  \item ``Compiler'' covers several genuinely different tools: \textbf{native/AOT compilers},
  \textbf{source-to-source transpilers}, \textbf{cross-compilers}, \textbf{binary translators},
  \textbf{decompilers}, and \textbf{bootstrapping/self-hosting compilers} -- know the
  source$\to$target shape of each.
  \item Compiler techniques (automata, grammars, data-flow analysis) show up far outside
  traditional compilers -- in IDEs, static analyzers, database engines, and ML frameworks -- which
  is the real reason this course is core, not elective.
  \item DB \S1.5's five application areas (language implementation, architecture optimization,
  architecture design, program translation, productivity tools) have since been joined by
  ML compilation, WebAssembly, and shared infrastructure like LLVM.
\end{itemize}
\end{takeaways}

\section*{References for This Lecture}
\begin{itemize}
  \item [DB] Aho, Lam, Sethi, Ullman. \textit{Compilers: Principles, Techniques, and Tools}, 2nd ed., \S1.1, \S1.2, \S1.4, \S1.5.
  \item [ET] Cooper \& Torczon. \textit{Engineering a Compiler}, 2nd ed., \S1.1.
  \item [AP] Appel. \textit{Modern Compiler Implementation in C/Java/ML}, Ch.\,1.
\end{itemize}

\end{document}
